% Copyright (c) 2026 Geovana Neves. Licensed under CC BY 4.0 (https://creativecommons.org/licenses/by/4.0/). See LICENSE-AND-AUTHORSHIP.md.
%
% fsi/text/07-examples.tex
%
% The numbers are what the three examples print when run against the tag's
% sources (python examples/<name>.py, PYTHONPATH at the tag's src/, the [fsi]
% extra installed), on 2026-09-28. Every blade is synthetic by the package's
% policy. A release that changes an example changes these numbers: run it
% again and copy them.

\section{Worked examples}

\begin{frame}[fragile]{Example 1: calculated and supplied inputs agree}
\begin{columns}[T]
\begin{column}{0.54\textwidth}
\code{examples/workspace\_fsi\_calibration.py}, offline, no solver:
\begin{enumerate}\scriptsize
  \item Writes \code{inputs/fsi/f\_calculated.toml} from the package template
        (the titanium strip), with the file factor 1.05 on EI.
  \item Resolves it and reads the \emph{unscaled} base distributions.
  \item Writes \code{f\_supplied.toml} with those base distributions and the
        same file factor 1.05. Reusing the effective values here would apply
        the calibration a second time.
  \item Resolves both with \code{FSI\_BENDING\_STIFFNESS\_N\_M2\_FACTOR: 1.15}.
  \item Asserts every distribution equal across the two modes, effective EI
        equal to base EI times 1.15, and the origin \code{matrix}.
\end{enumerate}
\begin{lstlisting}[style=ftsbase, basicstyle=\ttfamily\tiny]
python workspace_fsi_calibration.py /absolute/path/to/new-workspace
\end{lstlisting}
\end{column}
\begin{column}{0.42\textwidth}
{\ftstablesize\renewcommand{\arraystretch}{1.2}
\begin{tabular}{@{}lr@{}}
\toprule
\textbf{Printed} & \textbf{Value} \\
\midrule
base EI, each station & 24.277 N\,m\(^2\) \\
effective EI, each station & 27.919 N\,m\(^2\) \\
file factor, matrix factor & 1.05, 1.15 \\
\code{factor\_origin} & \code{matrix} \\
\code{mode\_equivalence} & verified \\
\code{solver\_executed} & \code{false} \\
\bottomrule
\end{tabular}}
\begin{ftsevidence}[Check it by hand]
\scriptsize \(I = b h^3/12 = 0.04 \times 0.004^3 / 12 = 2.133\times10^{-10}\) m\(^4\);
\(E I = 113.8\) GPa \(\times\, I = 24.28\) N\,m\(^2\);
\(\times\, 1.15 = 27.92\), not \(\times\, 1.05 \times 1.15\).
\end{ftsevidence}
{\scriptsize Choose a new destination: the example refuses to overwrite its
two inputs.\par}
\end{column}
\end{columns}
\fsisource{\code{examples/workspace\_fsi\_calibration.py}, run at the tag;
  \code{workspace/fsi\_setup.py}, \code{FSI\_TEMPLATE}}
\end{frame}

\begin{frame}{Example 2: one section, and its torsion constant}
\begin{columns}[T]
\begin{column}{0.48\textwidth}
\code{examples/fsi\_solid\_blade\_properties.py} builds a synthetic blade:
seven stations from 0.15 to 0.75 m, chord 80 to 40 mm, pitch 35 to 11 deg,
NACA 2415 at the root thinning to NACA 2409 at the tip, written to a section
file and read back, the file's sha256 into the provenance.
\par\vspace{2mm}
{\tiny\renewcommand{\arraystretch}{1.15}
\begin{tabular}{@{}lr@{}}
\toprule
\textbf{Root section, NACA 2415} & \\
\midrule
area & 654.144 mm\(^2\) \\
centroid from the pitch axis & \(-13.422\) mm chordwise \\
\(I\) flap, \(I\) chord & 5512.2, 227784.7 mm\(^4\) \\
principal angle & \(-0.100\) deg \\
\bottomrule
\end{tabular}}
\end{column}
\begin{column}{0.48\textwidth}
{\ftstablesize\renewcommand{\arraystretch}{1.15}
\begin{tabular}{@{}lr@{}}
\toprule
\textbf{Grid, cells across} & \(J\) [mm\(^4\)] \\
\midrule
16 & 20926.02 \\
32 & 20956.81 \\
64, the default & 20962.52 \\
thin-section estimate & 21721.14 \\
\bottomrule
\end{tabular}}
\par\vspace{1.5mm}
\begin{itemize}\scriptsize
  \item Doubling the cells divides the error by about four: at 64 cells about
        a third of the change from 32 to 64 is left.
  \item The thin-section formula is 3.6\,\% high here: a cross-check, never
        the value.
\end{itemize}
\end{column}
\end{columns}
\fsisource{\code{examples/fsi\_solid\_blade\_properties.py}, run at the tag}
\end{frame}

\begin{frame}{Example 2: the whole blade, in titanium}
\relax % a body opening with a brace group would be read as the frame subtitle
{\ftstablesize\renewcommand{\arraystretch}{1.1}
\begin{tabular*}{\linewidth}{@{\extracolsep{\fill}}rrrrrrrr@{}}
\toprule
\(r\) [m] & \(\mu\) [kg/m] & EI [N\,m\(^2\)] & GJ [N\,m\(^2\)] & \(I_1\) [kg\,m] & \(I_2\) [kg\,m] & \(e_c\) [mm] & thin\(/J\) \\
\midrule
0.15 & 2.8979 & 627.29 & 922.35 & 1.009e-03 & 2.442e-05 & \(-13.42\) & 1.036 \\
0.25 & 2.2727 & 361.01 & 531.76 & 6.650e-04 & 1.405e-05 & \(-12.30\) & 1.032 \\
0.35 & 1.7441 & 198.04 & 291.97 & 4.217e-04 & 7.708e-06 & \(-11.19\) & 1.028 \\
0.45 & 1.3040 & 102.60 & 151.24 & 2.554e-04 & 3.993e-06 & \(-10.07\) & 1.024 \\
0.55 & 0.9445 & 49.58 & 72.98 & 1.462e-04 & 1.930e-06 & \(-8.95\) & 1.020 \\
0.65 & 0.6574 & 21.98 & 32.24 & 7.789e-05 & 8.553e-07 & \(-7.83\) & 1.017 \\
0.75 & 0.4347 & 8.73 & 12.72 & 3.784e-05 & 3.395e-07 & \(-6.71\) & 1.014 \\
\bottomrule
\end{tabular*}}
\par\vspace{2mm}
\begin{columns}[T]
\begin{column}{0.56\textwidth}
\begin{itemize}\scriptsize
  \item Ti-6Al-4V grade 5, annealed: \(\rho\) 4430 kg/m\(^3\), \(E\) 113.8 GPa,
        \(G\) 44.0 GPa, \(\nu\) 0.342, database version 1.
  \item The configuration round-trips through \code{config.json} unchanged,
        with its provenance and its \code{config\_sha256}.
  \item \(e_c\): the centroid ahead of or behind the pitch axis, which the
        elastic axis is assumed to follow.
\end{itemize}
\end{column}
\begin{column}{0.40\textwidth}
\begin{ftsevidence}[At rest, from the beam]
\small first flap: 30.6 Hz\\
first torsion: 601.9 Hz
\end{ftsevidence}
\end{column}
\end{columns}
\fsisource{\code{examples/fsi\_solid\_blade\_properties.py}, sections 3 and 4,
  run at the tag}
\end{frame}

\begin{frame}{Example 3: the Campbell diagram of a rotating blade}
\begin{columns}[T]
\begin{column}{0.58\textwidth}
\begin{center}
\begin{adjustbox}{max width=\linewidth, max totalheight=0.70\textheight}
\begin{tikzpicture}[x=0.085cm, y=0.0118cm]
  % axes
  \draw[->, SoftGray] (0,0) -- (98,0);
  \node[font=\tiny, anchor=north] at (48,-30) {rotor speed \(\Omega\) [rad/s]};
  \draw[->, SoftGray] (0,0) -- (0,430);
  \node[font=\tiny, anchor=south, rotate=90] at (-14,215) {natural frequency \(\omega_n\) [rad/s]};
  \foreach \x in {0,15,30,45,60,75,90}
    \draw[SoftGray] (\x,0) -- (\x,-6) node[below, font=\tiny, text=black] {\x};
  \foreach \y in {0,100,200,300,400}
    \draw[SoftGray] (0,\y) -- (-1.5,\y) node[left, font=\tiny, text=black] {\y};
  % legend, in the empty upper left
  \draw[thick, CourseBlue] (5,405) -- (13,405);
  \fill[CourseBlue] (9,405) circle (1.2pt);
  \node[font=\tiny, anchor=west] at (14,405) {first flap};
  \draw[thick, CourseRed] (5,378) -- (13,378);
  \fill[CourseRed] (8.2,375.5) rectangle (9.8,380.5);
  \node[font=\tiny, anchor=west] at (14,378) {first torsion};
  \draw[dotted, SoftGray, thick] (5,351) -- (13,351);
  \node[font=\tiny, anchor=west] at (14,351) {\(n\)P excitation, \(n\,\Omega\)};
  % excitation rays nP
  \foreach \n in {1,2,3,4} {
    \draw[dotted, SoftGray, thick] (0,0) -- (90,{\n*90});
    \node[font=\tiny, text=SoftGray, right] at (90,{\n*90}) {\n P};
  }
  % first flap
  \draw[thick, CourseBlue] plot[mark=*, mark size=1.2pt] coordinates
    {(0,27.17) (15,31.70) (30,42.37) (45,55.55) (60,69.69) (75,84.23) (90,98.95)};
  % first torsion
  \draw[thick, CourseRed] plot[mark=square*, mark size=1.2pt] coordinates
    {(0,310.90) (15,311.25) (30,312.29) (45,314.02) (60,316.42) (75,319.48) (90,323.18)};
\end{tikzpicture}
\end{adjustbox}
\end{center}
\end{column}
\begin{column}{0.38\textwidth}
{\scriptsize A uniform synthetic blade, 1 m, 15 stations, swept from 0 to 90
rad/s. At every speed: static solve under its own centrifugal loads, then the
modes with the tension's geometric stiffness.\par}
\vspace{1.5mm}
{\ftstablesize\renewcommand{\arraystretch}{1.15}
\begin{tabular}{@{}lrr@{}}
\toprule
\textbf{Southwell fit} & \textbf{flap} & \textbf{torsion} \\
\midrule
\(\omega_0\) [rad/s] & 27.96 & 310.90 \\
\(S\) & 1.118 & 0.961 \\
\(r^2\) & 0.99986 & 1.00000 \\
\bottomrule
\end{tabular}}
\par\vspace{1.5mm}
{\scriptsize \(\omega_n^2 = \omega_0^2 + S\,\Omega^2\). Flap \(S\) sits in
the expected 1.1 to 1.3; torsion near \((I_1 - I_2)/(I_1 + I_2) = 0.961\).
A crossing of a track and a ray flags a potential resonance [10, 11].\par}
\end{column}
\end{columns}
\fsisource{\code{examples/fsi\_campbell\_diagram.py}, run at the tag;
  FSI tutorial, ``centrifugal''}
\end{frame}
